\documentclass[10pt,a4paper]{amsart}
\usepackage[margin=19mm]{geometry}
\usepackage{amsmath,amssymb,mathtools}
\usepackage[scr=rsfs,cal=euler]{mathalfa}
\usepackage{xfrac}
\usepackage[unicode,hidelinks,bookmarks=false]{hyperref}
\newcommand{\ie}{i.e.\ }
\newcommand{\cf}{cf.\ }
\newcommand{\resp}{respectively\ }
\newcommand{\Subsection}{\subsection}
\providecommand{\nobreakdash}{\nobreak\hbox{-}}
\setlength{\emergencystretch}{2em}
\multlinegap0pt
\usepackage{bm}
\usepackage{bbm}
\usepackage{dsfont}
\usepackage{xspace}
\usepackage{cancel}
\usepackage{mathtools}
\usepackage[shortlabels]{enumitem}
\usepackage{amsthm}
\makeatletter
\@ifundefined{proposition}{\newtheorem{proposition}{Proposition}}{}
\@ifundefined{Notation}{\newtheorem{Notation}[proposition]{Notation}}{}
\makeatother
\DeclareMathOperator{\Aut}{Aut}
\DeclareMathOperator{\BMix}{BMix}
\DeclareMathOperator{\HH}{H}
\DeclareMathOperator{\InnAut}{InnAut}
\DeclareMathOperator{\Mix}{Mix}
\DeclareMathOperator{\ZMix}{ZMix}
\newcommand{\CC}{\mathbb{C}}
\newcommand{\PP}{\mathbb{P}}
\newcommand{\W}{\mathcal{W}}
\newcommand{\ZZ}{\mathbb{Z}}
\newcommand{\case}[2]{\vspace{2mm} \textbf{Case #1:} #2 \vspace{2mm}}
\newcommand{\nonzero}{\setminus \{0\}}
\title[Central extensions, derivations, and automorphisms of semi-d]{Central extensions, derivations, and automorphisms of semi-direct sums of the Witt algebra with its intermediate series modules}
\author{Lucas Buzaglo, Girish S. Vishwa}
\date{Source: arXiv:2407.14809v2 (CC BY 4.0)}
\begin{document}
\maketitle

\noindent\emph{Source and licence.} Central extensions, derivations, and automorphisms of semi-direct sums of the Witt algebra with its intermediate series modules. Version arXiv:2407.14809v2 (CC BY 4.0), distributed under the Creative Commons Attribution 4.0 International licence, as recorded in the arXiv licence metadata for this exact version and shown on the abstract page https://arxiv.org/abs/2407.14809v2. Full rights notice: licenses/arxiv-2407.14809-CC-BY-4.0.txt.\par\smallskip

\noindent\textit{Editorial scope.} Counted unit: the Proposition whose source label is \texttt{prop:outer automorphisms} in the article (source0.tex lines 1743--1746, source label \texttt{prop:outer automorphisms}), delivered together with the article's own complete original proof of it (source0.tex lines 1747--1853, one \texttt{proof} environment of 107 lines). No mathematics is written, repaired or re-derived: statement and proof are the article's own text line for line. Required context retained uncounted: ntt:abelian cocycle B (lines 1012--1014); prop:W\_A mixing cocycles (lines 1099--1112); eq:mixing cocycle A (lines 1126--1128); ntt:mixing cocycles B (lines 1199--1203); prop:W\_B mixing cocycles (lines 1205--1211); eq:simple mixing cocycle B (lines 1241--1243); eq:outer automorphism general form (lines 1499--1501); prop:automorphism nice form (lines 1688--1692). Every definition, display and label of those lines is retained unchanged. Omitted from this package and not redistributed: 1--1011, 1015--1098, 1113--1125, 1129--1198, 1204--1204, 1212--1240, 1244--1498, 1502--1687, 1693--1742, 1854--2478. Nothing inside a retained excerpt is omitted or altered: every retained body is delivered exactly as the article's lines are written, so no omission receipt of any kind accompanies this package. Citations: the retained excerpts carry no citation command at all, so no bibliography entry of the article is transcribed and no reference is redistributed. Reference adaptation: none was needed and none was made; every reference inside the delivered unit resolves inside it, so no unresolved reference or citation is produced. Printed slips kept literal: no printed word, symbol, hypothesis or number of the counted statement or of its proof was altered. Layout and class: the article's own class is \texttt{amsart} with packages including fontenc, enumitem, multirow, xcolor, chngcntr, adjustbox, comment, import, xifthen, appendix, etoolbox, geometry, amsfonts, amsmath; the wrapper is the lane's pinned \texttt{amsart} editorial wrapper at 10pt on A4, which restates the theorem-like environments the retained text needs and adds the article's own macro definitions that the retained text needs (\texttt{\textbackslash Aut}, \texttt{\textbackslash BMix}, \texttt{\textbackslash CC}, \texttt{\textbackslash HH}, \texttt{\textbackslash InnAut}, \texttt{\textbackslash Mix}, \texttt{\textbackslash PP}, \texttt{\textbackslash W}, \texttt{\textbackslash ZMix}, \texttt{\textbackslash ZZ}, \texttt{\textbackslash case}, \texttt{\textbackslash nonzero}), with extra wrapper packages (bm, bbm, dsfont, xspace, cancel, mathtools, enumitem). The delivered numbering is the unit's own. Assets: no retained span contains a figure environment, a graphics-inclusion command, a PDF-page inclusion, a table or a TikZ picture; no image file is copied into this package, the package declares no asset dependency and no image file is redistributed with it. Licence: version arXiv:2407.14809v2 of this article is distributed under the Creative Commons Attribution 4.0 International licence, as recorded in the arXiv licence metadata for this exact version and shown on the abstract page https://arxiv.org/abs/2407.14809v2. A full rights notice is delivered at licenses/arxiv-2407.14809-CC-BY-4.0.txt. Characters: every retained excerpt is pure ASCII, so the delivered engine, XeLaTeX with its default Latin Modern text font, reports no missing character.
\begin{Notation}\label{ntt:abelian cocycle B}
    Define $\iota \colon \ZZ \to \CC$ by $\iota(n) = n$ for all $n \in \ZZ$.
\end{Notation}

\begin{proposition}\label{prop:W_A mixing cocycles}
    Let $\lambda \in \PP^1$. Then $\ZMix_A(0) = \CC \beta_0 \oplus \CC \iota$, where $\iota$ is as in Notation \ref{ntt:abelian cocycle B}, and $\ZMix_A(\lambda) = \CC \beta_\lambda$ if $\lambda \neq 0$. In particular, if $\Omega$ is a mixing cocycle on $\W_A(\lambda)$ of degree zero, then for $\lambda \neq 0$ we have, up to multiplication by a scalar,
    $$\Omega(L_n,A_m) = \begin{cases}
        (\lambda + 1)\delta_{n+m}^0, &\text{if } n = 0 \text{ and } \lambda \neq \infty, \\
        \delta_{n+m}^0, &\text{otherwise},
    \end{cases}$$
    while for $\lambda = 0$, there exist $a,b \in \CC$ such that
    $$\Omega(L_n,A_m) = (an + b)\delta_{n+m}^0.$$
    Furthermore, $\BMix_A(\lambda) = 0$, so
    $$\dim(\HH_{\Mix}^2(\W_A(\lambda))) = \begin{cases}
        2, &\text{if } \lambda = 0,\\
        1, &\text{if } \lambda \neq 0.
    \end{cases}$$
\end{proposition}

    \begin{equation}\label{eq:mixing cocycle A}
        (-n + m(m + 1)\lambda \delta_{n+m}^0)\beta(n) - (-m + n(n + 1)\lambda \delta_{n+m}^0)\beta(m) + (n - m)\beta(n + m) = 0.
    \end{equation}

\begin{Notation}\label{ntt:mixing cocycles B}
    Let $\lambda \in \PP^1$. Define $\eta_\lambda, \gamma_1, \gamma_2 \colon \ZZ \to \CC$ as follows:
    $$\eta_\lambda(n) = n + n(n + 1)\lambda, \quad \gamma_1(n) = n, \quad \gamma_2(n) = n^2$$
    for all $n \in \ZZ$, where $n(n + 1)\lambda$ is understood to be $n^2$ if $\lambda = \infty$. In other words, $\eta_\lambda = (\lambda + 1) \gamma_1 + \lambda \gamma_2$ if $\lambda \neq \infty$ and $\eta_\infty = \gamma_1 + \gamma_2$.
\end{Notation}

\begin{proposition}\label{prop:W_B mixing cocycles}
    Let $\lambda \in \PP^1$. Then $\ZMix_B(\lambda) = \CC \gamma_1 \oplus \CC \gamma_2$, while $\BMix_B(\lambda) = \CC \eta_\lambda$, where $\gamma_1,\gamma_2$ and $\eta_\lambda$ are as in Notation \ref{ntt:mixing cocycles B}. In other words, if $\Omega$ is a mixing cocycle of degree 0 on $\W_B(\lambda)$, then
    $$\Omega(L_n,B_m) = (an^2 + bn)\delta_{n+m}^0$$
    for some $a,b \in \CC$. Furthermore, if $\lambda \neq \infty$, then $\Omega$ is a coboundary if and only if there exists $c \in \CC$ such that $a = c\lambda$ and $b = c(\lambda + 1)$, while if $\lambda = \infty$, then $\Omega$ is a coboundary if and only if $a = b$.
    Consequently,
    $$\dim(\HH_{\Mix}^2(\W_B(\lambda))) = \dim(\ZMix_B(\lambda)) - \dim(\BMix_B(\lambda)) = 1.$$
\end{proposition}

\begin{equation}\label{eq:simple mixing cocycle B}
    (n + m)(\beta(m) - \beta(n)) + (n - m)\beta(n + m) = 0.
\end{equation}

    \begin{equation}\label{eq:outer automorphism general form}
        \Phi_{k;a,b;\alpha,\xi}^{X(\lambda)} = \tau_\lambda^k \circ \varphi_a^{X(\lambda)} \circ \psi_b^{X(\lambda)} \circ \sigma_\alpha \circ \mu_\xi,
    \end{equation}

\begin{proposition}\label{prop:automorphism nice form}
    Let $X \in \{A,B\}$, $\lambda \in \PP^1$, and $\sigma \in \Aut(\W_X(\lambda))$. Then there exists $\rho \in \InnAut(\W_X(\lambda))$ such that, for all $n \in \ZZ$,
    $$(\rho \circ \sigma)(L_n) = \varepsilon \alpha^n L_{\varepsilon n} + \gamma_n X_{\varepsilon n}, \quad (\rho \circ \sigma)(X_n) = \xi_n X_{\varepsilon n},$$
    for some $\varepsilon \in \{\pm 1\}$, $\alpha, \xi_n \in \CC^\times$, $\gamma_n \in \CC$. Furthermore, if $X = A$, then $\gamma_0 = 0$.
\end{proposition}

\begin{proposition}\label{prop:outer automorphisms}
    Let $X \in \{A,B\}$, $\lambda \in \PP^1$, and $\sigma \in \Aut(\W_X(\lambda))$. Then there exists $\rho \in \InnAut(\W_X(\lambda))$ such that $\rho \circ \sigma = \Phi_{k;a,b;\alpha,\xi}^{X(\lambda)}$ for some $k \in \ZZ/2\ZZ$, $a,b \in \CC$, and $\alpha, \xi \in \CC^\times$, where $\Phi_{k;a,b;\alpha,\xi}^{X(\lambda)}$ is defined in \eqref{eq:outer automorphism general form}.

\end{proposition}

\begin{proof}
    Let $\sigma \in \Aut(\W_X(\lambda))$. By Proposition \ref{prop:automorphism nice form}, possibly by replacing $\sigma$ with $\rho \circ \sigma$ for some $\rho \in \InnAut(\W_X(\lambda))$, we may assume that
    \begin{equation}\label{eq:sigma nice form}
        \sigma(L_n) = \alpha^n (\varepsilon L_{\varepsilon n} + \gamma_n X_{\varepsilon n}), \quad \sigma(X_n) = \alpha^n \xi_n X_{\varepsilon n},
    \end{equation}
    for some $\varepsilon \in \{\pm 1\}$, $\alpha, \xi_n \in \CC^\times$, $\gamma_n \in \CC$, with $\gamma_0 = 0$ if $X = A$.

    Since $[L_n,L_m] = (m - n)L_{n+m}$, we have $[\sigma(L_n),\sigma(L_m)] = (m - n)\sigma(L_{n+m})$. Therefore, applying \eqref{eq:sigma nice form} and comparing coefficients of $X_{\varepsilon n}$, we deduce that
    \begin{equation}\label{eq:automorphism gamma}
        \varepsilon(\omega_{X(\lambda)}(\varepsilon n, \varepsilon m)\gamma_m - \omega_{X(\lambda)}(\varepsilon m, \varepsilon n)\gamma_n) = (m - n)\gamma_{n+m}
    \end{equation}
    for all $n,m \in \ZZ$. Similarly, we have $[\sigma(L_n),\sigma(X_m)] = \omega_{X(\lambda)}(n,m)\sigma(X_{n+m})$, so we can apply \eqref{eq:sigma nice form} and simplify to get
    \begin{equation}\label{eq:automorphism xi}
        \varepsilon \omega_{X(\lambda)}(\varepsilon n, \varepsilon m)\xi_m = \omega_{X(\lambda)}(n,m)\xi_{n+m}
    \end{equation}
    for all $n,m \in \ZZ$. Furthermore, any bijective linear map $\sigma$ defined as in \eqref{eq:sigma nice form} which satisfies \eqref{eq:automorphism gamma} and \eqref{eq:automorphism xi} is an automorphism of $\W_X(\lambda)$.
    
    We now consider the cases $X = A$ and $X = B$ separately.

    \case{1}{$X = A$.}
    
    Recall that in this case, we have $\gamma_0 = 0$. It follows from \eqref{eq:automorphism gamma} and the definition of $\omega_{A(\lambda)}$ that
    \begin{equation}\label{eq:automorphism gamma A}
        (n + m + n(\varepsilon n + 1)\lambda \delta_m^0)\gamma_m - (n + m + m(\varepsilon m + 1)\lambda \delta_n^0)\gamma_n = (m - n)\gamma_{n+m}.
    \end{equation}
    Since $\gamma_0 = 0$, we have $\delta_n^0 \gamma_n = 0$ for all $n \in \ZZ$. Therefore, we can simplify \eqref{eq:automorphism gamma A} as follows:
    \begin{equation}\label{eq:automorphism gamma A simple}
        (n + m)(\gamma_m - \gamma_n) = (m - n)\gamma_{n+m}
    \end{equation}
    for all $n \in \ZZ$. Note that \eqref{eq:automorphism gamma A simple} has the same form as \eqref{eq:simple mixing cocycle B}, which we already solved in the proof of Proposition \ref{prop:W_B mixing cocycles}: we have $\gamma_n = a n^2 + b n$ for some $a,b \in \CC$ (in fact, $a = \frac{1}{2}\gamma_2 - \gamma_1$ and $b = 2\gamma_1 - \frac{1}{2}\gamma_2$).

    We now compute $\xi_n$. By \eqref{eq:automorphism xi}, we have
    \begin{equation}\label{eq:automorphism xi A}
        (n + m + n(\varepsilon n + 1)\lambda \delta_m^0)\xi_m = (n + m + n(n + 1)\lambda \delta_m^0)\xi_{n+m}
    \end{equation}
    for all $n,m \in \ZZ$. It follows from \eqref{eq:automorphism xi A} that $\xi_m = \xi_{n+m}$ for all $m \in \ZZ \nonzero$ and $n \in \ZZ \setminus \{-m\}$. Letting $\xi = \xi_1$, we therefore conclude that $\xi_n = \xi$ for all $n \in \ZZ \nonzero$.

    To compute $\xi_0$, we substitute $n \in \ZZ \nonzero$ and $m = 0$ into \eqref{eq:automorphism xi A} to get
    \begin{equation}\label{eq:xi0 A}
        (n + n(\varepsilon n + 1)\lambda)\xi_0 = (n + n(n + 1)\lambda)\xi
    \end{equation}
    for all $n \in \ZZ$. If $\varepsilon = 1$, then \eqref{eq:xi0 A} implies that $\xi_0 = \xi$. It follows that $\sigma = \varphi_a^{A(\lambda)} \circ \psi_b^{A(\lambda)} \circ \sigma_\alpha \circ \mu_\xi$ (in other words, $\sigma = \Phi^{A(\lambda)}_{0,a,b,\alpha,\xi}$).
    
    Now assume that $\varepsilon = -1$. In this case, if $\lambda \neq \infty$, \eqref{eq:xi0 A} becomes
    $$(\xi + \xi_0)\lambda n + (\xi - \xi_0)(\lambda + 1) = 0$$
    for all $n \in \ZZ \nonzero$. Therefore, $(\xi + \xi_0)\lambda = 0$ and $(\xi - \xi_0)(\lambda + 1) = 0$. Similarly, if $\lambda = \infty$, then \eqref{eq:xi0 A} gives $(\xi + \xi_0)n + \xi - \xi_0 = 0$ for all $n \in \ZZ \nonzero$, from which it follows that $\xi_0 = \xi = 0$. We conclude that, when $\varepsilon = -1$, there are three possibilities:
    \begin{enumerate}[(a)]
        \item $\lambda = 0$ and $\xi_0 = \xi$;
        \item $\lambda = -1$ and $\xi_0 = -\xi$;
        \item $\lambda \in \PP^1 \setminus \{0,-1\}$ and $\xi_0 = \xi = 0$.
    \end{enumerate}
    However, $\sigma$ is an automorphism, so we cannot have $\xi = 0$. Therefore, the case $\varepsilon = -1$ is not possible if $\lambda \in \PP^1 \setminus \{0,-1\}$. Therefore, in the two cases where $\varepsilon = -1$ is possible ($\lambda = 0$ or $\lambda = -1$), we get $\sigma = \tau_\lambda \circ \varphi_a^{A(\lambda)} \circ \psi_b^{A(\lambda)} \circ \sigma_\alpha \circ \mu_\xi$ (in other words, $\sigma = \Phi_{1;a,b;\alpha,\xi}^{A(\lambda)}$). This concludes the proof when $X = A$.

    \case{2}{$X = B$.}

    Let $\gamma_n' = \gamma_{\varepsilon n}$ and $k = \varepsilon n$, $\ell = \varepsilon m$. Then \eqref{eq:automorphism gamma} becomes
    \begin{equation}\label{eq:automorphism gamma' B}
        \omega_{B(\lambda)}(k,\ell)\gamma_\ell' - \omega_{B(\lambda)}(\ell,k)\gamma_k' = (\ell - k)\gamma_{k+\ell}',
    \end{equation}
    for all $k,\ell \in \ZZ$. By \eqref{eq:automorphism gamma' B} and the definition of $\omega_{B(\lambda)}$, we have
    \begin{equation}\label{eq:automorphism gamma B simple}
        (\ell - k(k + 1)\lambda \delta_{k+\ell}^0)\gamma_\ell' - (k - \ell(\ell + 1)\lambda \delta_{k+\ell}^0)\gamma_k' = (\ell - k)\gamma_{k+\ell}'.
    \end{equation}
    Note that \eqref{eq:automorphism gamma B simple} has the same form as \eqref{eq:mixing cocycle A}, so we already know the solutions, thanks to the proof of Proposition \ref{prop:W_A mixing cocycles}: if $\lambda = 0$, then $\gamma_n' = a'n + b$ for some $a',b \in \CC$, while for $\lambda \notin \{0,\infty\}$, we get
    $$\gamma_n' = \begin{cases}
        (\lambda + 1)c, &\text{if } n = 0; \\
        c, &\text{if } n \neq 0.
    \end{cases}$$
    for some $c \in \CC$. For $\lambda = \infty$, we get $\gamma_n' = c$ for all $n \in \ZZ$. Therefore, letting $a = \varepsilon a'$, we conclude that $\gamma_n = an + b$ if $\lambda = 0$, while
    $$\gamma_n = \begin{cases}
        (\lambda + 1)c, &\text{if } n = 0; \\
        c, &\text{if } n \neq 0.
    \end{cases}$$
    if $\lambda \notin \{0,\infty\}$, and $\gamma_n = c$ for all $n \in \ZZ$ if $\lambda = \infty$.

    We finish by computing $\xi_n$. By \eqref{eq:automorphism xi}, we get
    \begin{equation}\label{eq:automorphism xi B}
        (m - n(\varepsilon n + 1)\lambda \delta_{n+m}^0)\xi_m = (m - n(n + 1)\lambda \delta_{n+m}^0)\xi_{n+m}
    \end{equation}
    for all $n,m \in \ZZ$. Similarly to before, letting $\xi = \xi_1$, \eqref{eq:automorphism xi B} immediately implies that $\xi_n = \xi$ for all $n \in \ZZ \nonzero$.
    
    Finally, in order to compute $\xi_0$, we substitute $m \in \ZZ \nonzero$ and $n = -m$ into \eqref{eq:automorphism xi B}:
    \begin{equation}\label{eq:xi0 B}
        (m - m(\varepsilon m - 1)\lambda)\xi = (m - m(m - 1)\lambda)\xi_0
    \end{equation}
    for all $m \in \ZZ \nonzero$. If $\varepsilon = 1$, then it follows from \eqref{eq:xi0 B} that $\xi_0 = \xi$. Hence,
    $$\sigma = \begin{cases}
        \varphi_a^{B(0)} \circ \psi_b^{B(0)} \circ \sigma_\alpha \circ \mu_\xi, &\text{if } \lambda = 0, \\
        \varphi_c^{B(\lambda)} \circ \sigma_\alpha \circ \mu_\xi, &\text{if } \lambda \neq 0.
    \end{cases}$$
    In other words, $\sigma = \Phi_{0;a,b,\alpha,\xi}^{B(0)}$ if $\lambda = 0$, while $\sigma = \Phi_{0;c,0;\alpha,\xi}^{B(\lambda)}$ if $\lambda \neq 0$.
    
    Therefore, it remains to consider $\varepsilon = -1$. In this case, if $\lambda \neq \infty$, \eqref{eq:xi0 B} becomes
    $$(\xi + \xi_0)\lambda m + (\xi - \xi_0)(\lambda + 1) = 0$$
    for all $m \in \ZZ \nonzero$. For $\lambda = \infty$, we instead get $(\xi + \xi_0)m + \xi - \xi_0 = 0$ for all $m \in \ZZ \nonzero$. Proceeding as we did in Case 1, we conclude that there are three possibilities when $\varepsilon = -1$:
    \begin{enumerate}[(a)]
        \item $\lambda = 0$ and $\xi_0 = \xi$;
        \item $\lambda = -1$ and $\xi_0 = -\xi$;
        \item $\lambda \in \PP^1 \setminus \{0,-1\}$ and $\xi_0 = \xi = 0$.
    \end{enumerate}
    It follows that, as before, the case $\varepsilon = -1$ is not possible if $\lambda \in \PP^1 \setminus \{0,-1\}$. Therefore, in the cases where $\varepsilon = -1$ is possible ($\lambda = 0$ or $\lambda = -1$), we get
    $$\sigma = \begin{cases}
        \tau_0 \circ \varphi_a^{B(0)} \circ \psi_b^{B(0)} \circ \sigma_\alpha \circ \mu_\xi, &\text{if } \lambda = 0, \\
        \tau_{-1} \circ \varphi_c^{B(-1)} \circ \sigma_\alpha \circ \mu_\xi, &\text{if } \lambda = -1.
    \end{cases}$$
    In other words, $\sigma = \Phi_{1;a,b;\alpha,\xi}^{B(0)}$ if $\lambda = 0$, while $\sigma = \Phi_{1;c,0;\alpha,\xi}^{B(-1)}$ if $\lambda = -1$, which concludes the proof.
\end{proof}

\end{document}
